\section{Strong clipped energy under total-variation convergence of signed data}
\label{sec:r11-signed-data-strong-clipping}

Retain the original arbitrary degree-$(n-1,n-1)$ background,
the affinoid neighborhood and the ordinary relatively compact
part domain of the signed diffuse Dirichlet theorem.  All
measures below are finite signed energy-diffuse measures on
$F$, and all potentials use their canonical energy-quasicontinuous
representatives.  The norm of $H$ is $\mathcal E^{1/2}$.

\begin{theorem}[Strong clipping and smooth-curvature stability]
\label{thm:r11-signed-data-strong-clipping}
Suppose $\|\mu_j-\mu\|_{\rm TV}\to0$, and write
$u_j=u_{\mu_j}$, $u=u_\mu$ for the unique intrinsic
renormalized solutions of $D(u_jS)=\mu_j$, $D(uS)=\mu$.
Then, for every fixed $k>0$,
\[
                    \|T_ku_j-T_ku\|_H\longrightarrow0.
 \tag{DS1}
\]
In particular, for any fixed $k,l>0$, the finite original
energy measures converge in global total variation:
\[
 \Gamma_S(T_ku_j,T_lu_j)
               \longrightarrow\Gamma_S(T_ku,T_lu).
 \tag{DS2}
\]
If $\Phi\in C^2(\mathbf R)$ has bounded first derivative
and compactly supported second derivative, then
\[
           D(\Phi(u_j)S)\longrightarrow D(\Phi(u)S)
                       \quad\hbox{in global total variation}.
 \tag{DS3}
\]
If, more restrictively, $\Phi'$ has compact support, then
$\Phi(u_j)-\Phi(0)\to\Phi(u)-\Phi(0)$ strongly in $H$.
For this latter class both original first-order currents
$d'(\Phi(u_j)S)$ and $d''(\Phi(u_j)S)$ converge in local
coefficient total variation, against every original compact
complementary smooth test and every actual moment refinement.
The zero-order multipliers $\Phi(u_j)S$ converge in that same
local coefficient sense under the hypotheses of (DS3).
\end{theorem}

\noindent\textbf{Status: \texttt{DIRECT\_COROLLARY}} of
the completed signed renormalized theorem, especially the exact
band-energy identity (SD8), and the proved finite-datum capacity
stability.  This conclusion is not used to prove (SD8).

\begin{proof}
The finite-datum estimates give $u_j\to u$ in energy capacity
and in $L^1(\lambda)$.  The clipped functions converge in
$L^2(\lambda)$: they are bounded by $k$, and $\lambda(F)<\infty$
because the part domain is relatively compact in the neighborhood
where $\lambda$ is Radon.  Also
$\mathcal E(T_ku_j)\leq k\|\mu_j\|_{\rm TV}$, so these
vectors are bounded in $H$.  Every weakly convergent subsequence
therefore has the unique $L^2$ representative $T_ku$.
Consequently the whole sequence converges weakly in $H$.

Take $a=0,b=k$ in (SD8).  The exact norm identity is
\[
 \mathcal E(T_ku_j)=-\int_F T_k\widetilde u_j\,d\mu_j,
 \qquad
 \mathcal E(T_ku)=-\int_F T_k\widetilde u\,d\mu.
 \tag{DS4}
\]
The difference between the first datum integral and the second
is bounded in absolute value by
\[
 k\|\mu_j-\mu\|_{\rm TV}
       +\int_F|T_k\widetilde u_j-T_k\widetilde u|\,d|\mu|.
 \tag{DS5}
\]
The second term tends to zero by energy-capacity convergence,
finite diffuseness of $|\mu|$, and boundedness of the clippings.
Thus the norms in (DS4) converge.  Weak convergence and convergence
of norms in the Hilbert space prove (DS1).  In particular the
argument does not infer strong clipping merely from a uniform
energy bound and an $L^1$ limit; its decisive input is the exact
identity already proved in the signed renormalized theorem.

The original closed gradient module gives, for $f_j,g_j,f,g\in H$,
\[
 \|\Gamma_S(f_j,g_j)-\Gamma_S(f,g)\|_{\rm TV}
  \leq \|f_j-f\|_H\|g_j\|_H
                         +\|f\|_H\|g_j-g\|_H.
 \tag{DS6}
\]
This is measure Cauchy--Schwarz after adding and subtracting
$\Gamma_S(f,g_j)$.  Applying (DS1) proves (DS2).

For (DS3), choose $M$ larger than the support radius of $\Phi''$.
The all-original chain identity reads
\[
 D(\Phi(u_j)S)=\Phi'(\widetilde u_j)\mu_j
       +\Phi''(\widetilde u_j)\Gamma_S(T_Mu_j).
 \tag{DS7}
\]
The first term converges in variation: its changing-datum error
is at most $\|\Phi'\|_\infty\|\mu_j-\mu\|_{\rm TV}$,
and its fixed-datum multiplier error tends to zero by finite
diffuseness and capacity convergence.  For the second term use
(DS2) and then bounded convergence in measure against the fixed
finite diffuse measure $\Gamma_S(T_Mu)$.  This gives its variation
convergence as well and proves (DS3).  The zero-order multiplier
claim follows from the Lipschitz bound for $\Phi$, $L^1(\lambda)$
stability and the faithful original coefficient domination.

If $\Phi'$ has compact support, its normalized composition is
a smooth function of one sufficiently large clipping.  The
closed gradient formula and (DS1) give strong convergence:
the changing-gradient term is bounded by
$\|\Phi'\|_\infty\|T_Mu_j-T_Mu\|_H$, and the
changing scalar factor tends to zero in the fixed finite
energy measure.  Its $L^2$ representative converges by the
same bounded composition argument.  Finally, on a fixed
original compact stage, write any complementary one-form
test as a finite sum of bounded smooth coefficients times
differentials of original local moment functions.  Multiply
those moment functions by original compact cutoffs equal to
one on the test support.  The coefficient pairing of a
first-order derivative with each such term is an original
mixed gradient measure; (DS6), or its one-variable
Cauchy--Schwarz bound, controls its variation by a fixed
constant times the $H$-norm of the changing potential.
This proves local coefficient-TV convergence of both first-order
currents.  The argument starts with an arbitrary original test
and is valid after any further actual moment refinement.
It is not a claim of one global variation norm for all
lower-degree coefficients simultaneously.
\end{proof}

\begin{proposition}[Hard clipping curvatures need not be TV continuous in the datum]
\label{prop:r11-hard-clipping-source-TV-failure}
For a fixed $k>0$, the implication
\[
 \mu_j\longrightarrow\mu\text{ in total variation}
 \quad\Longrightarrow\quad
 D(T_ku_{\mu_j}S)\longrightarrow D(T_ku_\mu S)
                   \text{ in total variation}
 \tag{DS8}
\]
is \texttt{FALSE}, even for positive smooth original data on an
original annulus.  This does not contradict (DS1)--(DS3) or
the global hard-curvature limit as $k\to\infty$ in (SD11).
\end{proposition}

\noindent\textbf{Status: \texttt{DIRECT\_COROLLARY}} of the
frozen all-original one-dimensional torus calibration, with the
following explicit computation.  The rejected inference (DS8)
is recorded as \texttt{FALSE}.

\begin{proof}
Let $V=\{x\in\mathbf G_{m,K}^{\rm an}:-1<t(x)<1\}$,
where $t=\log|T|$, and take $S=[V]$.  It extends to a larger
ordinary annulus inside one original affinoid.  Write
$\sigma:(-1,1)\to V$ for the original Gauss section.  Put
\[
 u=t^2-1,\qquad u_j=c_j u,\qquad c_j=1+j^{-1},
 \qquad \mu_j=2c_j\sigma_*(dt),\quad\mu=2\sigma_*(dt).
 \tag{DS9}
\]
These potentials belong to the part space $H$: smooth real
compact cutoffs approximate their zero-boundary real Sobolev
profiles in gradient energy, and the exact original calibration
passes that approximation to the full original form.  Their
all-original curvature measures are the displayed data.
They are therefore the canonical solutions by the finite-energy
or signed renormalized uniqueness theorem.  In particular these
finite data are energy diffuse and
$\|\mu_j-\mu\|_{\rm TV}=4(c_j-1)\to0$.

Fix $0<k<1$, and put
\[
                  a_j=\sqrt{1-k/c_j},\qquad a=\sqrt{1-k}.
\]
Since the potentials are nonpositive, $T_ku_j=\max(u_j,-k)$.
Its real slope is zero on $(-a_j,a_j)$ and $2c_jt$ outside.
The two positive slope jumps have size $2c_ja_j$.  The
original calibrated measure is consequently
\[
 D(T_ku_jS)
  =\sigma_*\left(2c_j\mathbf1_{\{|t|>a_j\}}\,dt
                +2c_ja_j(\delta_{-a_j}+\delta_{a_j})\right).
 \tag{DS10}
\]
The same formula with $c_j=1,a_j=a$ gives $D(T_kuS)$.
The Gauss section is injective, $a_j>a$, and all four atoms
in the difference are distinct.  The diffuse parts converge
in total variation, while the total variation of the atomic
difference tends to
\[
                         4c_ja_j+4a\longrightarrow8a>0.
 \tag{DS11}
\]
Thus the total variation in (DS8) does not tend to zero.
For an arbitrary prescribed $k>0$, replace $u$ by
$A(t^2-1)$ with $A>k$, and replace $u_j$ by $c_j u$.
The same calculation has $a_j=\sqrt{1-k/(Ac_j)}$ and
limiting atomic variation $8A\sqrt{1-k/A}>0$.
Each equality is an original Radon measure equality from the
frozen calibration, valid on every original scalar test and
every actual moment refinement.  Strong convergence of the
clipped gradients still holds by (DS1); it is the moving
hard-level curvature atoms, not their energy measures, that
prevent the stronger asserted continuity.
\end{proof}

\paragraph{Dependency boundary.}
The new stability statement is downstream of the signed
renormalized chain and exact band identity.  It does not close
their initial proof by a circular norm argument.  The smooth
scalar curvature limit (DS3) has its full derivative hypotheses;
the explicit failure (DS8) prevents replacing those hypotheses
by an arbitrary hard clamp.  Unweighted energy of an unbounded
solution may still fail to be locally Radon, and no ambient
classical potential or ordinary pointwise boundary value is
asserted.
